\documentclass[UTF8,12pt,a4paper]{ctexart}

\usepackage{amsmath}
\usepackage{cases}
\usepackage{cite}
\usepackage{graphicx}
\usepackage{enumerate}
\usepackage{algorithm}
\usepackage{caption} %\caption*需要
\usepackage[noend]{algpseudocode} %algorithmicx
\makeatletter
\renewcommand{\fnum@algorithm}{\fname@algorithm}
\makeatother
\renewcommand{\algorithmicrequire}{\textbf{Input:}}
\renewcommand{\algorithmicensure}{\textbf{Output:}}
\usepackage[margin=1in]{geometry}
\geometry{a4paper}
\usepackage{fancyhdr}
\pagestyle{fancy}
\fancyhf{}
\usepackage{xcolor}
\usepackage{listings}
\lstset{
	basicstyle=\ttfamily,                                % 传统上，代码打字机字体好些
	breaklines,
	tabsize=4,
	extendedchars=false,
	columns=fixed,
	%numbers=left,                                        % 在左侧显示行号
	numbers=none,
	frame=none,                                          % 不显示背景边框
	backgroundcolor=\color[RGB]{245,245,244},            % 设定背景颜色
	keywordstyle=\color[RGB]{40,40,255},                 % 设定关键字颜色
	numberstyle=\footnotesize\color{darkgray},           % 设定行号格式
	commentstyle=\it\color[RGB]{0,96,96},                % 设置代码注释的格式
	stringstyle=\rmfamily\slshape\color[RGB]{128,0,0},   % 设置字符串格式
	showstringspaces=false,                              % 不显示字符串中的空格
	xleftmargin=1em,
	%xrightmargin=1em,
	%aboveskip=-0.5em
	language=c++,                                        % 设置语言
}

% title依据是作业还是实验报告，以及是第几次作业/实验报告，
% 比如第1次分治算法作业，则命名为"算法原理 HW1 分治算法";
% 第1次分治算法实验报告，则命名为"算法原理 Lab1 分治算法"
\title{算法原理 HW5 作业}
\author{
	姓名：\underline{冯峻}~~~~~~
	学号：\underline{523031910148}~~~~~~}
\date{\today}
\pagenumbering{arabic}

\begin{document}

% 如下两行请按实际情况修改
\fancyhead[L]{冯峻}
\fancyhead[C]{随机算法}
\fancyfoot[C]{\thepage}

\maketitle
%\tableofcontents
%\newpage

\section*{作业题目}
\begin{enumerate}[1.]
	\item 有一个函数 $F:\{0,1,\ldots,n-1\} \to \{0,1,\ldots,m-1\}$，且 $F((x+y)\bmod n) = (F(x) + F(y))\bmod m$ 对 $\forall x, y \in \{0,1,\ldots,n-1\}$ 成立，其中 $\bmod$ 为求模操作。设 $F(x)$ 存储在一个数组中，数组下标表示自变量的值，数组元素的值表示函数值；由于某种意外，数组中 $1/5$ 的函数值被恶意篡改。试设计一个随机算法使其对 $\forall z \in \{0,1,\ldots,n-1\}$，算法能够以大于 $1/2$ 的概率计算出正确的 $F(z)$。如果算法运行 3 次，应该返回什么样的值？此时算法得到正确 $F(z)$ 的概率有什么变化？
\end{enumerate}

\section*{问题分析}

本题的关键在于利用函数的加法同态性质：$F((x+y)\bmod n) = (F(x) + F(y))\bmod m$。

对于任意 $z \in \{0,1,\ldots,n-1\}$，我们可以随机选择 $x \in \{0,1,\ldots,n-1\}$，计算 $y = (z-x) \bmod n$，则有 $(x+y) \bmod n = z$。根据函数的性质，应该满足：
$$F(z) = (F(x) + F(y)) \bmod m$$

由于只有 $1/5$ 的值被篡改，因此有 $4/5$ 的值是正确的。当我们随机选择 $x$ 和 $y$ 时：
\begin{itemize}
	\item $F(x)$ 和 $F(y)$ 都正确的概率为 $(4/5) \times (4/5) = 16/25 = 0.64$
	\item 至少一个错误的概率为 $1 - 16/25 = 9/25 = 0.36$
\end{itemize}

因此，单次随机验证得到正确 $F(z)$ 的概率为 $16/25 > 1/2$，满足题目要求。

\section*{算法设计}

\subsection{单次随机验证算法（Monte Carlo 算法）}

基本思想是随机选择 $x$，利用函数的加法性质计算候选值。算法伪代码如下：

\begin{algorithm}[H]
	\caption{COMPUTE-F($F, n, m, z$)}
	\begin{algorithmic}[1]
		\Require 函数数组 $F$，定义域大小 $n$，值域大小 $m$，查询位置 $z$
		\Ensure $F(z)$ 的估计值
		\State 随机选择 $x \in \{0,1,\ldots,n-1\}$
		\State $y \gets (z - x + n) \bmod n$
		\State $candidate \gets (F[x] + F[y]) \bmod m$
		\State \Return{$candidate$}
	\end{algorithmic}
\end{algorithm}

\textbf{正确性分析：}算法返回正确值的概率为 $P = (4/5)^2 = 16/25 = 0.64 > 1/2$。

\subsection{多次验证与投票机制}

为了提高算法的准确性，可以进行 $k$ 次随机验证，并采用多数投票（Majority Voting）策略选择出现次数最多的候选值。

\begin{algorithm}[H]
	\caption{COMPUTE-F-VOTING($F, n, m, z, k$)}
	\begin{algorithmic}[1]
		\Require 函数数组 $F$，定义域大小 $n$，值域大小 $m$，查询位置 $z$，验证次数 $k$
		\Ensure $F(z)$ 的估计值
		\State 初始化投票字典 $votes \gets \emptyset$
		\For{$i = 1$ to $k$}
			\State 随机选择 $x \in \{0,1,\ldots,n-1\}$
			\State $y \gets (z - x + n) \bmod n$
			\State $candidate \gets (F[x] + F[y]) \bmod m$
			\State $votes[candidate] \gets votes[candidate] + 1$
		\EndFor
		\State \Return{获得最多票数的 $candidate$}
	\end{algorithmic}
\end{algorithm}

\subsection{算法运行 3 次的策略}

当算法运行 3 次时，应该\textbf{返回出现次数最多的值（众数）}。这是因为：

\begin{itemize}
	\item 单次运行得到正确结果的概率为 $p = 16/25 = 0.64$
	\item 运行 3 次，至少 2 次正确（从而众数正确）的概率为：
	$$P_{correct} = \binom{3}{2}p^2(1-p) + \binom{3}{3}p^3$$
	$$= 3 \times (16/25)^2 \times (9/25) + (16/25)^3$$
	$$= 3 \times 0.4096 \times 0.36 + 0.262144$$
	$$= 0.442368 + 0.262144 = 0.704512 \approx 0.70$$
\end{itemize}

因此，运行 3 次并返回众数，正确概率从 $0.64$ 提升到约 $0.70$，概率有所提高。

\section*{算法实现}

\subsection{C++ 代码实现}

算法的 C++ 代码如下：

\begin{lstlisting}[language={C++}, basicstyle=\ttfamily\small]
#include <iostream>
#include <vector>
#include <map>
#include <cstdlib>
#include <ctime>
using namespace std;

// 单次随机验证算法
int singleRandomCheck(vector<int>& F, int n, int m, int z) {
    int x = rand() % n;              // 随机选择x
    int y = (z - x + n) % n;         // 计算y
    return (F[x] + F[y]) % m;        // 返回候选值
}

// 运行k次并投票
int computeWithVoting(vector<int>& F, int n, int m, int z, int k) {
    map<int, int> votes;  // 存储投票结果
    
    for (int i = 0; i < k; i++) {
        int candidate = singleRandomCheck(F, n, m, z);
        votes[candidate]++;
    }
    
    // 返回获得最多票数的值
    int maxVotes = 0, result = F[z];
    for (auto& p : votes) {
        if (p.second > maxVotes) {
            maxVotes = p.second;
            result = p.first;
        }
    }
    return result;
}

// 运行3次算法并返回众数
int threeTimesAlgorithm(vector<int>& F, int n, int m, int z) {
    return computeWithVoting(F, n, m, z, 3);
}

int main() {
    srand(time(0));
    int n = 10, m = 10;
    
    // 示例：F(x) = (2*x) mod m
    vector<int> originalF(n);
    for (int i = 0; i < n; i++)
        originalF[i] = (2 * i) % m;
    
    // 篡改1/5的值
    vector<int> corruptedF = originalF;
    for (int i = 0; i < n/5; i++) {
        int pos = rand() % n;
        corruptedF[pos] = rand() % m;
    }
    
    // 测试算法
    int z = 5;
    cout << "真实值: " << originalF[z] << endl;
    cout << "单次验证: " << singleRandomCheck(corruptedF, n, m, z) << endl;
    cout << "3次验证: " << threeTimesAlgorithm(corruptedF, n, m, z) << endl;
    
    return 0;
}
\end{lstlisting}

\section*{总结}

本题设计的随机算法是一个 Monte Carlo 算法，其主要特点是：
\begin{itemize}
	\item 利用函数的加法同态性质进行验证
	\item 单次运行的正确概率为 $16/25 = 0.64 > 1/2$
	\item 通过多次运行和投票机制可以进一步提高准确率
	\item 运行 3 次时采用众数投票，正确概率提升至约 $0.70$
	\item 算法时间复杂度为 $O(k)$，其中 $k$ 是验证次数
\end{itemize}



\end{document}
